%%%%%%%%%%%%%%%%%%%%%%%%%%%%%%%%%%%%%%%%%%%%%%%%%%%%%%%%%%%%%%%%%%%%%%%
% Thema: Poster
% Autor: Musterman
% Datum: 12.05
%%%%%%%%%%%%%%%%%%%%%%%%%%%%%%%%%%%%%%%%%%%%%%%%%%%%%%%%%%%%%%%%%%%%%%%
%
\documentclass[a0,portrait]{a0poster}
%\documentclass[portrait,a0,posterdraft]{a0poster}
%
%%%%%%%%%%%%%%%%%%%%%%%%%%%%%%%%%%%%%%%%%%%%%%%%%%%%%%%%%%%%%%%%%%%%%%%
% Paket zum Erzeugen der Spalten
\usepackage{multicol}

%%%%%%%%%%%%%%%%%%%%%%%%%%%%%%%%%%%%%%%%%%%%%%%%%%%%%%%%%%%%%%%%%%%%%%%
% Eingabekodierung
\usepackage[ansinew]{inputenc}
\usepackage[T1]{fontenc}
\usepackage{ae} %wenn Ulaute nicht gehen

\usepackage[ngerman]{babel}
%\begin{�bler Trick um die Orinal�berschrift zu unterdr�cken}
\addto{\captionsngerman}{\renewcommand*{\refname}{}} %Literatur�berschfit unterdr�cken
%\ende{�bler Trick}

%%%%%%%%%%%%%%%%%%%%%%%%%%%%%%%%%%%%%%%%%%%%%%%%%%%%%%%%%%%%%%%%%%%%%%%
% f�r Farbe
\usepackage{color} 
\definecolor{darkgreen}{rgb}{0,0.5,0} 
\definecolor{darkblue}{rgb}{0,0,0.5}

%%%%%%%%%%%%%%%%%%%%%%%%%%%%%%%%%%%%%%%%%%%%%%%%%%%%%%%%%%%%%%%%%%%%%%%
% Mathepaket
\usepackage{amsmath} 

%%%%%%%%%%%%%%%%%%%%%%%%%%%%%%%%%%%%%%%%%%%%%%%%%%%%%%%%%%%%%%%%%%%%%%%
% Ist f�r Boxen (wie das hier benutzte Ovalbox) zust�ndig
\usepackage{fancybox}

%%%%%%%%%%%%%%%%%%%%%%%%%%%%%%%%%%%%%%%%%%%%%%%%%%%%%%%%%%%%%%%%%%%%%%%
% Graphikpaket, erm�glicht png, jpg und pdf bilder
\usepackage{graphicx}

%%%%%%%%%%%%%%%%%%%%%%%%%%%%%%%%%%%%%%%%%%%%%%%%%%%%%%%%%%%%%%%%%%%%%%%
% Seiteneinstellungen
\renewcommand\baselinestretch{1.35}
\parskip=0.5\baselineskip

\parindent0mm %Einr�cktiefe der ersten Zeile eines Absatzes
\topmargin-28pt
\marginparwidth0mm

%R�nder rechts/links
\oddsidemargin-13pt
\evensidemargin-13pt
\textwidth785mm
%\textheight1140mm

%%%%%%%%%%%%%%%%%%%%%%%%%%%%%%%%%%%%%%%%%%%%%%%%%%%%%%%%%%%%%%%%%%%%%%%
% Eigene Definitionen zur Erleichterung des Satzes
\newcommand{\spaltenbreite}{25}   % Spaltenbreite f�r Bilder
\newcommand{\bildbreite}{25cm}    % Einheitliche Bildbreite

%%%%%%%%%%%%%%%%%%%%%%%%%%%%%%%%%%%%%%%%%%%%%%%%%%%%%%%%%%%%%%%%%%%%%%%
% Box- und Spalteneinstellungen
\setlength{\fboxrule}{3.25mm} %Definiert die Linienst�rke f�r nachfolgende fbox- und framebox-Befehle
\setlength{\fboxsep}{5mm} %Abstand zwischen Rahmen und Text bei den /fbox und /framebox Befehlen.
\setlength{\columnsep}{15mm}     %Spaltenabstand
\setlength{\columnseprule}{1pt}  %Balken zwischen Spalten {0pt}->keine Balken

%%%%%%%%%%%%%%%%%%%%%%%%%%%%%%%%%%%%%%%%%%%%%%%%%%%%%%%%%%%%%%%%%%%%%%%
% Einheitsl�nge f�r picture-Umgebungen
%\setlength{\unitlength}{1.0cm}
\unitlength1cm

%%%%%%%%%%%%%%%%%%%%%%%%%%%%%%%%%%%%%%%%%%%%%%%%%%%%%%%%%%%%%%%%%%%%%%%
% Grafikpfad, hier liegen alle Bilder
%\graphicspath{{images/}}

%%%%%%%%%%%%%%%%%%%%%%%%%%%%%%%%%%%%%%%%%%%%%%%%%%%%%%%%%%%%%%%%%%%%%%%
% Zaehler fuer lineale. Sie werden gebraucht, wenn das Linealmacro
% included wird
\newcounter{skalax}
\newcounter{skalay}


\begin{document}

%%%%%%%%%%%%%%%%%%%%%%%%%%%%%%%%%%%%%%%%%%%%%%%%%%%%%%%%%%%%%%%%%%%%%%%
% Kopf, hier funktioniert alles mit 'put'
%%%%%%%%%%%%%%%%%%%%%%%%%%%%%%%%%%%%%%%%%%%%%%%%%%%%%%%%%%%%%%%%%%%%%%%
%\parbox{\textwidth}
%{
  \begin{center}
    \vspace*{0.006\textheight}
   % Titel und Autor
    {\huge \textbf{�bersicht �ber Schaltbare Spiegel}\\}%[0.03\textheight]
    \vspace*{0.03\textheight}
    {\large \textbf{Gruppe 15: Nima Astan Varzeghan, Marc Alexander Depping, Andr�s Felipe Fonseca Richter,\\Dimitrios Maximilian Kalodikis, Luka Lolovic, Juan Nicolas Pardo Martin}}
  \end{center}
  % Logos einbinden (per Hand plazieren)
  \begin{picture}(0,0)
    \put(70.0 , 4.6){\includegraphics[height=80mm]{logo}}
    \put(71.8 , 2.8){\LARGE \textsf{\textbf{TU Berlin}}}
    \put(0.0  , 4.3){\includegraphics[height=83mm]{logo}}
    \put(0.0  , 3.1){\large \textsf{\textbf{Funktionswerkstoffe der Elektrotechnik}}}
    \put(0.0  , 1.75){\large \textsf{\textbf{Prof. Dr.-Ing. Ronald Plath}}}
    \put(0.0  , 0.4){\large \textsf{\textbf{Prof. Dr. rer. nat. Bernd Szyszka}}}
  \end{picture}
%}

%%%%%%%%%%%%%%%%%%%%%%%%%%%%%%%%%%%%%%%%%%%%%%%%%%%%%%%%%%%%%%%%%%%%%%%
% Einstellung des Eckenradius' der ovalen Boxen (\Ovalbox)
\cornersize*{5mm}


\linethickness{0.1mm}
\setlength{\fboxrule}{2.25mm}

%%%%%%%%%%%%%%%%%%%%%%%%%%%%%%%%%%%%%%%%%%%%%%%%%%%%%%%%%%%%%%%%%%%%%%% 
% neuer Kasten
\Ovalbox
{
  \parbox{\textwidth}{
    \begin{multicols}{3}
      
            \begin{center} \textbf{\Large Abschnitt A} \end{center}
Zkjkl \cite{blackstock} kljl�dksjf iouepoitueoiut �lak�sdlkfj  uopweuroiwe eri iweuro wiueoweor 
ioruwerilkvklvmkl meio difji ifjeoij  ijfiowej. Zkjkl kljl�dksjf iouepoitueoiut 
�lak�sdlkfj  uopweuroiwe eri iweuro wiueoweor ioruwerilkvklvmkl meio difji ifjeoij
  ijfiowej. Zkjkl kljl�dksjf iouepoitueoiut �lak�sdlkfj  uopweuroiwe eri iweuro
 wiueoweor ioruwerilkvklvmkl meio difji ifjeoij  ijfiowej. Zkjkl kljl�dksjf 
iouepoitueoiut �lak�sdlkfj  

Zlak�sdlkfj \cite{drp} uopweuroiwe eri iweuro wiueoweor ioruwerilkvklvmkl meio difji ifjeoij  
kwjheorw jsku3w48oi9 jhk9 r98 jkhkf fdgre.
A lkdso dlksloiio  lsdo lwflkweklf offeowlkd io kdjfljweoie sdkfowpe
kwjheorw jsku3w48oi9 jhk9 r98 jkhkf \cite{drp} fdgre.

            \begin{center} \textbf{\Large Abschnitt B} \end{center}
Zkjkl \cite{blackstock} kljlödksjf iouepoitueoiut ölakösdlkfj  uopweuroiwe eri iweuro wiueoweor 
ioruwerilkvklvmkl meio difji ifjeoij  ijfiowej. Zkjkl kljlödksjf iouepoitueoiut 
ölakösdlkfj  uopweuroiwe eri iweuro wiueoweor ioruwerilkvklvmkl meio difji ifjeoij
  ijfiowej. Zkjkl kljlödksjf iouepoitueoiut ölakösdlkfj  uopweuroiwe eri iweuro
 wiueoweor ioruwerilkvklvmkl meio difji ifjeoij  ijfiowej. Zkjkl kljlödksjf 
iouepoitueoiut ölakösdlkfj  

uopweuroiwe eri iweuro wiueoweor ioruwerilkvklvmkl meio difji ifjeoij  ijfiowej. 
Zkjkl kljlödksjf iouepoitueoiut ölakösdlkfj  uopweuroiwe eri iweuro wiueoweor 
ioruwerilkvklvmkl meio difji ifjeoij  ijfiowej. Zkjkl kljlödksjf iouepoitueoiut 
ölakösdlkfj  uopweuroiwe eri iweuro wiueoweor ioruwerilkvklvmkl meio difji ifjeoij  

ijfiowej. Zkjkl kljlödksjf iouepoitueoiut ölakösdlkfj  uopweuroiwe eri iweuro 
wiueoweor ioruwerilkvklvmkl meio difji ifjeoij  ijfiowej. Zkjkl kljlödksjf 
iouepoitueoiut ölakösdlkfj  kl kljlödksjf iouepoit.
A lkdso dlksloiio  lsdo lwflkweklf offeowlkd io iojoiojoi kdjfljweoie sdkfowpe
kslduoiwue9874 930875  hjfojd 3489 ölkj ölkj0ppooi4 44 kldjfljowi lskdjfoi 
kwjheorw jsku3w48oi9 jhk9 r98 jkhkf fdgre.
A lkdso dlksloiio  lsdo lwflkweklf offeowlkd io kdjfljweoie sdkfowpe
kwjheorw jsku3w48oi9 jhk9 r98 jkhkf \cite{drp} fdgre.
\begin{center}
\begin{picture}(\spaltenbreite, 12.3)
  \put(4.8,1.6){\includegraphics[width=0.2\textwidth,height=0.13\textwidth]{logo}}
  \put(0.2,0.3){\parbox{720pt}{{\bf \small Abb. 1:} \small Vergleich der physikalischen und numerischen Wellenzahl}}
%%%%%%%%%%%%%%%%%%%%%%%%%%%%%%%%%%%%%%%%%%%%%%%%%%%%%%%%%%%%%%%%%%%%%%%%%%%%%%%%%%%%%%%%%%%%%%
%% rulers
% initialise counters
%\newcounter{skalax}
\setcounter{skalax}{-1}
%\newcounter{skalay}
\setcounter{skalay}{-1}
% put text
\multiput(0,-0.2)(1,0){31}{\stepcounter{skalax}\makebox(0,0)[t]{\scriptsize\arabic{skalax}}}
\multiput(-0.8,0)(0,1){31}{\stepcounter{skalay}\makebox(0,0)[l]{\scriptsize\arabic{skalay}}}
% put lines
\newcommand{\oldline}{\linethickness}
\linethickness{0.1mm}
\multiput(0,0)(0,1){30}{\line(1,0){30}} % fine grid
\multiput(0,0)(1,0){30}{\line(0,1){30}}
\linethickness{0.4mm}
\multiput(0,0)(0,5){7}{\line(1,0){30}}  % medium grid
\multiput(0,0)(5,0){7}{\line(0,1){30}}
\linethickness{0.8mm}
\multiput(0,0)(10,0){4}{\line(0,1){30}} % coarse grid
\multiput(0,0)(0,10){4}{\line(1,0){30}}
\linethickness{\oldline}
%%%%%%%%%%%%%%%%%%%%%%%%%%%%%%%%%%%%%%%%%%%%%%%%%%%%%%%%%%%%%%%%%%%%%%%%%%%%%%%%%%%%%%%%%%%%%

\end{picture}
\end{center}

            \begin{center} \textbf{\Large Abschnitt B} \end{center}
Zkjkl \cite{blackstock} kljlödksjf iouepoitueoiut ölakösdlkfj  uopweuroiwe eri iweuro wiueoweor 
ioruwerilkvklvmkl meio difji ifjeoij  ijfiowej. Zkjkl kljlödksjf iouepoitueoiut 
ölakösdlkfj  uopweuroiwe eri iweuro wiueoweor ioruwerilkvklvmkl meio difji ifjeoij
  ijfiowej. Zkjkl kljlödksjf iouepoitueoiut ölakösdlkfj  uopweuroiwe eri iweuro
 wiueoweor ioruwerilkvklvmkl meio difji ifjeoij  ijfiowej. Zkjkl kljlödksjf 
iouepoitueoiut ölakösdlkfj  

uopweuroiwe eri iweuro wiueoweor ioruwerilkvklvmkl meio difji ifjeoij  ijfiowej. 
Zkjkl kljlödksjf iouepoitueoiut ölakösdlkfj  uopweuroiwe eri iweuro wiueoweor 
ioruwerilkvklvmkl meio difji ifjeoij  ijfiowej. Zkjkl kljlödksjf iouepoitueoiut 
ölakösdlkfj  uopweuroiwe eri iweuro wiueoweor ioruwerilkvklvmkl meio difji ifjeoij  

ijfiowej. Zkjkl kljlödksjf iouepoitueoiut ölakösdlkfj  uopweuroiwe eri iweuro 
wiueoweor ioruwerilkvklvmkl meio difji ifjeoij  ijfiowej. Zkjkl kljlödksjf 
iouepoitueoiut ölakösdlkfj  kl kljlödksjf iouepoit.
A lkdso dlksloiio  lsdo lwflkweklf offeowlkd io iojoiojoi kdjfljweoie sdkfowpe
kslduoiwue9874 930875  hjfojd 3489 ölkj ölkj0ppooi4 44 kldjfljowi lskdjfoi 
kwjheorw jsku3w48oi9 jhk9 r98 jkhkf fdgre.
A lkdso dlksloiio  lsdo lwflkweklf offeowlkd io kdjfljweoie sdkfowpe
kwjheorw jsku3w48oi9 jhk9 r98 jkhkf \cite{drp} fdgre.
\begin{center}
\begin{picture}(\spaltenbreite, 12.3)
  \put(4.8,1.6){\includegraphics[width=0.2\textwidth,height=0.13\textwidth]{logo}}
  \put(0.2,0.3){\parbox{720pt}{{\bf \small Abb. 1:} \small Vergleich der physikalischen und numerischen Wellenzahl}}
%%%%%%%%%%%%%%%%%%%%%%%%%%%%%%%%%%%%%%%%%%%%%%%%%%%%%%%%%%%%%%%%%%%%%%%%%%%%%%%%%%%%%%%%%%%%%%
%% rulers
% initialise counters
%\newcounter{skalax}
\setcounter{skalax}{-1}
%\newcounter{skalay}
\setcounter{skalay}{-1}
% put text
\multiput(0,-0.2)(1,0){31}{\stepcounter{skalax}\makebox(0,0)[t]{\scriptsize\arabic{skalax}}}
\multiput(-0.8,0)(0,1){31}{\stepcounter{skalay}\makebox(0,0)[l]{\scriptsize\arabic{skalay}}}
% put lines
\newcommand{\oldline}{\linethickness}
\linethickness{0.1mm}
\multiput(0,0)(0,1){30}{\line(1,0){30}} % fine grid
\multiput(0,0)(1,0){30}{\line(0,1){30}}
\linethickness{0.4mm}
\multiput(0,0)(0,5){7}{\line(1,0){30}}  % medium grid
\multiput(0,0)(5,0){7}{\line(0,1){30}}
\linethickness{0.8mm}
\multiput(0,0)(10,0){4}{\line(0,1){30}} % coarse grid
\multiput(0,0)(0,10){4}{\line(1,0){30}}
\linethickness{\oldline}
%%%%%%%%%%%%%%%%%%%%%%%%%%%%%%%%%%%%%%%%%%%%%%%%%%%%%%%%%%%%%%%%%%%%%%%%%%%%%%%%%%%%%%%%%%%%%

\end{picture}
\end{center}

      
    \end{multicols}
  }
}

%%%%%%%%%%%%%%%%%%%%%%%%%%%%%%%%%%%%%%%%%%%%%%%%%%%%%%%%%%%%%%%%%%%%%%%
% neuer Kasten
\vspace*{0.002\textheight}
\Ovalbox
{
  \parbox{\textwidth}{
    \begin{multicols}{3}
      
            \begin{center} \textbf{\Large Abschnitt B} \end{center}
Zkjkl \cite{blackstock} kljlödksjf iouepoitueoiut ölakösdlkfj  uopweuroiwe eri iweuro wiueoweor 
ioruwerilkvklvmkl meio difji ifjeoij  ijfiowej. Zkjkl kljlödksjf iouepoitueoiut 
ölakösdlkfj  uopweuroiwe eri iweuro wiueoweor ioruwerilkvklvmkl meio difji ifjeoij
  ijfiowej. Zkjkl kljlödksjf iouepoitueoiut ölakösdlkfj  uopweuroiwe eri iweuro
 wiueoweor ioruwerilkvklvmkl meio difji ifjeoij  ijfiowej. Zkjkl kljlödksjf 
iouepoitueoiut ölakösdlkfj  

uopweuroiwe eri iweuro wiueoweor ioruwerilkvklvmkl meio difji ifjeoij  ijfiowej. 
Zkjkl kljlödksjf iouepoitueoiut ölakösdlkfj  uopweuroiwe eri iweuro wiueoweor 
ioruwerilkvklvmkl meio difji ifjeoij  ijfiowej. Zkjkl kljlödksjf iouepoitueoiut 
ölakösdlkfj  uopweuroiwe eri iweuro wiueoweor ioruwerilkvklvmkl meio difji ifjeoij  

ijfiowej. Zkjkl kljlödksjf iouepoitueoiut ölakösdlkfj  uopweuroiwe eri iweuro 
wiueoweor ioruwerilkvklvmkl meio difji ifjeoij  ijfiowej. Zkjkl kljlödksjf 
iouepoitueoiut ölakösdlkfj  kl kljlödksjf iouepoit.
A lkdso dlksloiio  lsdo lwflkweklf offeowlkd io iojoiojoi kdjfljweoie sdkfowpe
kslduoiwue9874 930875  hjfojd 3489 ölkj ölkj0ppooi4 44 kldjfljowi lskdjfoi 
kwjheorw jsku3w48oi9 jhk9 r98 jkhkf fdgre.
A lkdso dlksloiio  lsdo lwflkweklf offeowlkd io kdjfljweoie sdkfowpe
kwjheorw jsku3w48oi9 jhk9 r98 jkhkf \cite{drp} fdgre.
\begin{center}
\begin{picture}(\spaltenbreite, 12.3)
  \put(4.8,1.6){\includegraphics[width=0.2\textwidth,height=0.13\textwidth]{logo}}
  \put(0.2,0.3){\parbox{720pt}{{\bf \small Abb. 1:} \small Vergleich der physikalischen und numerischen Wellenzahl}}
%%%%%%%%%%%%%%%%%%%%%%%%%%%%%%%%%%%%%%%%%%%%%%%%%%%%%%%%%%%%%%%%%%%%%%%%%%%%%%%%%%%%%%%%%%%%%%
%% rulers
% initialise counters
%\newcounter{skalax}
\setcounter{skalax}{-1}
%\newcounter{skalay}
\setcounter{skalay}{-1}
% put text
\multiput(0,-0.2)(1,0){31}{\stepcounter{skalax}\makebox(0,0)[t]{\scriptsize\arabic{skalax}}}
\multiput(-0.8,0)(0,1){31}{\stepcounter{skalay}\makebox(0,0)[l]{\scriptsize\arabic{skalay}}}
% put lines
\newcommand{\oldline}{\linethickness}
\linethickness{0.1mm}
\multiput(0,0)(0,1){30}{\line(1,0){30}} % fine grid
\multiput(0,0)(1,0){30}{\line(0,1){30}}
\linethickness{0.4mm}
\multiput(0,0)(0,5){7}{\line(1,0){30}}  % medium grid
\multiput(0,0)(5,0){7}{\line(0,1){30}}
\linethickness{0.8mm}
\multiput(0,0)(10,0){4}{\line(0,1){30}} % coarse grid
\multiput(0,0)(0,10){4}{\line(1,0){30}}
\linethickness{\oldline}
%%%%%%%%%%%%%%%%%%%%%%%%%%%%%%%%%%%%%%%%%%%%%%%%%%%%%%%%%%%%%%%%%%%%%%%%%%%%%%%%%%%%%%%%%%%%%

\end{picture}
\end{center}

            \begin{center} \textbf{\Large Abschnitt B} \end{center}
Zkjkl \cite{blackstock} kljlödksjf iouepoitueoiut ölakösdlkfj  uopweuroiwe eri iweuro wiueoweor 
ioruwerilkvklvmkl meio difji ifjeoij  ijfiowej. Zkjkl kljlödksjf iouepoitueoiut 
ölakösdlkfj  uopweuroiwe eri iweuro wiueoweor ioruwerilkvklvmkl meio difji ifjeoij
  ijfiowej. Zkjkl kljlödksjf iouepoitueoiut ölakösdlkfj  uopweuroiwe eri iweuro
 wiueoweor ioruwerilkvklvmkl meio difji ifjeoij  ijfiowej. Zkjkl kljlödksjf 
iouepoitueoiut ölakösdlkfj  

uopweuroiwe eri iweuro wiueoweor ioruwerilkvklvmkl meio difji ifjeoij  ijfiowej. 
Zkjkl kljlödksjf iouepoitueoiut ölakösdlkfj  uopweuroiwe eri iweuro wiueoweor 
ioruwerilkvklvmkl meio difji ifjeoij  ijfiowej. Zkjkl kljlödksjf iouepoitueoiut 
ölakösdlkfj  uopweuroiwe eri iweuro wiueoweor ioruwerilkvklvmkl meio difji ifjeoij  

ijfiowej. Zkjkl kljlödksjf iouepoitueoiut ölakösdlkfj  uopweuroiwe eri iweuro 
wiueoweor ioruwerilkvklvmkl meio difji ifjeoij  ijfiowej. Zkjkl kljlödksjf 
iouepoitueoiut ölakösdlkfj  kl kljlödksjf iouepoit.
A lkdso dlksloiio  lsdo lwflkweklf offeowlkd io iojoiojoi kdjfljweoie sdkfowpe
kslduoiwue9874 930875  hjfojd 3489 ölkj ölkj0ppooi4 44 kldjfljowi lskdjfoi 
kwjheorw jsku3w48oi9 jhk9 r98 jkhkf fdgre.
A lkdso dlksloiio  lsdo lwflkweklf offeowlkd io kdjfljweoie sdkfowpe
kwjheorw jsku3w48oi9 jhk9 r98 jkhkf \cite{drp} fdgre.
\begin{center}
\begin{picture}(\spaltenbreite, 12.3)
  \put(4.8,1.6){\includegraphics[width=0.2\textwidth,height=0.13\textwidth]{logo}}
  \put(0.2,0.3){\parbox{720pt}{{\bf \small Abb. 1:} \small Vergleich der physikalischen und numerischen Wellenzahl}}
%%%%%%%%%%%%%%%%%%%%%%%%%%%%%%%%%%%%%%%%%%%%%%%%%%%%%%%%%%%%%%%%%%%%%%%%%%%%%%%%%%%%%%%%%%%%%%
%% rulers
% initialise counters
%\newcounter{skalax}
\setcounter{skalax}{-1}
%\newcounter{skalay}
\setcounter{skalay}{-1}
% put text
\multiput(0,-0.2)(1,0){31}{\stepcounter{skalax}\makebox(0,0)[t]{\scriptsize\arabic{skalax}}}
\multiput(-0.8,0)(0,1){31}{\stepcounter{skalay}\makebox(0,0)[l]{\scriptsize\arabic{skalay}}}
% put lines
\newcommand{\oldline}{\linethickness}
\linethickness{0.1mm}
\multiput(0,0)(0,1){30}{\line(1,0){30}} % fine grid
\multiput(0,0)(1,0){30}{\line(0,1){30}}
\linethickness{0.4mm}
\multiput(0,0)(0,5){7}{\line(1,0){30}}  % medium grid
\multiput(0,0)(5,0){7}{\line(0,1){30}}
\linethickness{0.8mm}
\multiput(0,0)(10,0){4}{\line(0,1){30}} % coarse grid
\multiput(0,0)(0,10){4}{\line(1,0){30}}
\linethickness{\oldline}
%%%%%%%%%%%%%%%%%%%%%%%%%%%%%%%%%%%%%%%%%%%%%%%%%%%%%%%%%%%%%%%%%%%%%%%%%%%%%%%%%%%%%%%%%%%%%

\end{picture}
\end{center}

      
    \end{multicols}
  }
}

%%%%%%%%%%%%%%%%%%%%%%%%%%%%%%%%%%%%%%%%%%%%%%%%%%%%%%%%%%%%%%%%%%%%%%%
% neuer Kasten
\vspace*{0.002\textheight}
\Ovalbox%\fbox
{
  \parbox{\textwidth}{
    \begin{multicols}{3}
      
            \begin{center} \textbf{\Large Abschnitt B} \end{center}
Zkjkl \cite{blackstock} kljlödksjf iouepoitueoiut ölakösdlkfj  uopweuroiwe eri iweuro wiueoweor 
ioruwerilkvklvmkl meio difji ifjeoij  ijfiowej. Zkjkl kljlödksjf iouepoitueoiut 
ölakösdlkfj  uopweuroiwe eri iweuro wiueoweor ioruwerilkvklvmkl meio difji ifjeoij
  ijfiowej. Zkjkl kljlödksjf iouepoitueoiut ölakösdlkfj  uopweuroiwe eri iweuro
 wiueoweor ioruwerilkvklvmkl meio difji ifjeoij  ijfiowej. Zkjkl kljlödksjf 
iouepoitueoiut ölakösdlkfj  

uopweuroiwe eri iweuro wiueoweor ioruwerilkvklvmkl meio difji ifjeoij  ijfiowej. 
Zkjkl kljlödksjf iouepoitueoiut ölakösdlkfj  uopweuroiwe eri iweuro wiueoweor 
ioruwerilkvklvmkl meio difji ifjeoij  ijfiowej. Zkjkl kljlödksjf iouepoitueoiut 
ölakösdlkfj  uopweuroiwe eri iweuro wiueoweor ioruwerilkvklvmkl meio difji ifjeoij  

ijfiowej. Zkjkl kljlödksjf iouepoitueoiut ölakösdlkfj  uopweuroiwe eri iweuro 
wiueoweor ioruwerilkvklvmkl meio difji ifjeoij  ijfiowej. Zkjkl kljlödksjf 
iouepoitueoiut ölakösdlkfj  kl kljlödksjf iouepoit.
A lkdso dlksloiio  lsdo lwflkweklf offeowlkd io iojoiojoi kdjfljweoie sdkfowpe
kslduoiwue9874 930875  hjfojd 3489 ölkj ölkj0ppooi4 44 kldjfljowi lskdjfoi 
kwjheorw jsku3w48oi9 jhk9 r98 jkhkf fdgre.
A lkdso dlksloiio  lsdo lwflkweklf offeowlkd io kdjfljweoie sdkfowpe
kwjheorw jsku3w48oi9 jhk9 r98 jkhkf \cite{drp} fdgre.
\begin{center}
\begin{picture}(\spaltenbreite, 12.3)
  \put(4.8,1.6){\includegraphics[width=0.2\textwidth,height=0.13\textwidth]{logo}}
  \put(0.2,0.3){\parbox{720pt}{{\bf \small Abb. 1:} \small Vergleich der physikalischen und numerischen Wellenzahl}}
%%%%%%%%%%%%%%%%%%%%%%%%%%%%%%%%%%%%%%%%%%%%%%%%%%%%%%%%%%%%%%%%%%%%%%%%%%%%%%%%%%%%%%%%%%%%%%
%% rulers
% initialise counters
%\newcounter{skalax}
\setcounter{skalax}{-1}
%\newcounter{skalay}
\setcounter{skalay}{-1}
% put text
\multiput(0,-0.2)(1,0){31}{\stepcounter{skalax}\makebox(0,0)[t]{\scriptsize\arabic{skalax}}}
\multiput(-0.8,0)(0,1){31}{\stepcounter{skalay}\makebox(0,0)[l]{\scriptsize\arabic{skalay}}}
% put lines
\newcommand{\oldline}{\linethickness}
\linethickness{0.1mm}
\multiput(0,0)(0,1){30}{\line(1,0){30}} % fine grid
\multiput(0,0)(1,0){30}{\line(0,1){30}}
\linethickness{0.4mm}
\multiput(0,0)(0,5){7}{\line(1,0){30}}  % medium grid
\multiput(0,0)(5,0){7}{\line(0,1){30}}
\linethickness{0.8mm}
\multiput(0,0)(10,0){4}{\line(0,1){30}} % coarse grid
\multiput(0,0)(0,10){4}{\line(1,0){30}}
\linethickness{\oldline}
%%%%%%%%%%%%%%%%%%%%%%%%%%%%%%%%%%%%%%%%%%%%%%%%%%%%%%%%%%%%%%%%%%%%%%%%%%%%%%%%%%%%%%%%%%%%%

\end{picture}
\end{center}

            \begin{center} \textbf{\Large Abschnitt A} \end{center}
Zkjkl \cite{blackstock} kljl�dksjf iouepoitueoiut �lak�sdlkfj  uopweuroiwe eri iweuro wiueoweor 
ioruwerilkvklvmkl meio difji ifjeoij  ijfiowej. Zkjkl kljl�dksjf iouepoitueoiut 
�lak�sdlkfj  uopweuroiwe eri iweuro wiueoweor ioruwerilkvklvmkl meio difji ifjeoij
  ijfiowej. Zkjkl kljl�dksjf iouepoitueoiut �lak�sdlkfj  uopweuroiwe eri iweuro
 wiueoweor ioruwerilkvklvmkl meio difji ifjeoij  ijfiowej. Zkjkl kljl�dksjf 
iouepoitueoiut �lak�sdlkfj  

Zlak�sdlkfj \cite{drp} uopweuroiwe eri iweuro wiueoweor ioruwerilkvklvmkl meio difji ifjeoij  
kwjheorw jsku3w48oi9 jhk9 r98 jkhkf fdgre.
A lkdso dlksloiio  lsdo lwflkweklf offeowlkd io kdjfljweoie sdkfowpe
kwjheorw jsku3w48oi9 jhk9 r98 jkhkf \cite{drp} fdgre.

      %%%%%%%%%%%%%%%%%%%%%%%%%%%%%%%%%%%%%%%%%%%%%%%%%%%%%%%%%%%%%%%%%%%%%%%
% Literaturverzeichnis

%\begin{�bler Trick um die Orinal�berschrift zu unterdr�cken}
\vspace*{1cm}
      \begin{center} \textbf{\Large Literatur} \end{center}
\vspace*{-4.5cm}
%\end{�bler Trick}

\renewcommand{\baselinestretch}{1.05}\small\normalsize
\begin{thebibliography}{9}
\bibitem{drp} \textsc{C. Tam \& J. Webb,} ``Dispersion-Relation-Preserving Finite Difference Schemes for Computational Acoustics'', \textsl{J. Comput. Phys.} 107, 1993.

\bibitem{lddrk} \textsc{F. Q. Hu, M. Y. Hussani \& J. L. Manthey,}  ``Low-Dissipation and -Dispersion Runge-Kutta Schemes for Computational Acoustics'',  \textsl{J. Comput. Phys.} 124(1), 1996.

\bibitem{shortwave} \textsc{C.K.W. Tam, J.C. Webb \& Z. Dong,} ``A study of the short wave components in computational acoustics'', \textsl{J. Comput. Acoust.} 1, 1993. 

\bibitem{hamilton} \textsc{Hamilton \& Blackstock}, \textit{Nonlinear Acoustics}, Academic Press, 1997.

\bibitem{blackstock} \textsc{Blackstock \& Hamilton}, \textit{Nonlinear Acoustics}, Academic Press, 1997.
\end{thebibliography}

      
    \end{multicols}
  }
}

\end{document}
